\needspace{10\baselineskip}
\section{进程线程模型}

本章的逻辑主线：为什么需要进程（多道程序设计带来并发环境，“程序”这个概念不够用了）$\to$ 怎样描述进程（状态、PCB、地址空间、队列）$\to$ 怎样控制进程（原语；UNIX 的 fork/exec/wait/exit）$\to$ 为什么又引入线程 $\to$ 线程的三种实现方式（用户级/核心级/混合）$\to$ 协程。答题时如果无从下手，就\textbf{从数据结构切入}：进程 $\to$ PCB；地址空间 $\to$ \cd{mm_struct} 和 VMA；线程 $\to$ TCB 和各自的栈。

\subsection{重点问题}

\Q{Q3-1.1 怎样理解“进程是对 CPU 的抽象”这句话？}
物理上只有一个（或几个）CPU，但系统中有许多程序要运行。OS 让\textbf{每个进程都以为自己独占了一个 CPU}：有自己的 PC、自己的寄存器、自己的执行流，从头到尾连续地执行。实际上 OS 通过\textbf{分时复用}——把 CPU 时间切成时间片，借助上下文切换在进程之间轮流分配——把一个物理 CPU 变成了多个\textbf{虚拟 CPU}，而进程就是这些虚拟 CPU 的载体。课件的说法：进程“将一个 CPU 变换成多个虚拟的 CPU”。

例（OSTEP \cd{cpu.c}）：在单核机器上同时运行 \cd{./cpu A & ./cpu B & ./cpu C &}，三个程序交替打印各自的字母。程序代码里没有任何“让出 CPU”的逻辑，每个程序都以为自己一直在运行，交替是 OS 在背后完成的。

这个抽象的实现 = \textbf{机制}（上下文切换：保存一个进程的寄存器、恢复另一个的）+ \textbf{策略}（调度算法：下一个运行谁）。同理，\textbf{地址空间是对内存的抽象}，\textbf{文件是对磁盘的抽象}。

\Q{Q3-1.2 何谓进程映像？进程有实体吗？在哪里？}
\textbf{进程映像}（process image）是进程在某一时刻的全部内容：\textbf{程序代码 + 数据 + 栈 + PCB}（Stallings 的定义；运行时还要加上 CPU 寄存器中的现场）。

进程是有实体的，它分布在以下几处：
\begin{itemize}
    \item \textbf{内存}：用户空间中的代码、数据、堆、栈，经页表映射到物理页框；内核空间中的 PCB 和该进程的内核栈。
    \item \textbf{磁盘}：还没有被调入内存的代码页仍在可执行文件里（按需调页）；被换出的页在交换区；进程被挂起时，整个映像都在磁盘上。
    \item \textbf{CPU 寄存器}：进程正在运行时，它的现场（PC、SP、通用寄存器、页表基址等）就在 CPU 寄存器中；不运行时这些值保存在 PCB 里。
\end{itemize}
xv6 中一个进程的实体：\cd{proc[NPROC]} 数组中的一个 \cd{struct proc}、一个内核栈 \cd{kstack}、一页 \cd{trapframe}、一张页表 \cd{pagetable} 以及它映射的那些物理页。

\Q{Q3-1.3 怎样描述一个进程？一个进程都有什么（组成要素）？}
用\textbf{进程控制块 PCB} 来描述。PCB 是 OS 用于管理进程的一个专门的数据结构，记录进程的各种属性，描述进程的动态变化过程。它与进程\textbf{一一对应}，是\textbf{系统感知进程存在的唯一标志}：创建进程就是创建 PCB，撤销进程就是回收 PCB。PCB 中的内容可以分为四类：
\begin{enumerate}
    \item \textbf{进程描述信息}：进程标识符 PID（唯一的整数）、进程名、用户标识符 UID、进程组关系（父进程、子进程）。
    \item \textbf{进程控制信息}：当前状态、优先级、代码执行入口地址、程序的磁盘地址（可执行文件位置）、运行统计信息（执行时间、页面调度次数）、进程间同步和通信（如阻塞原因、消息队列指针）、进程的队列指针。
    \item \textbf{所拥有的资源和使用情况}：虚拟地址空间（页表）、打开文件列表等。
    \item \textbf{CPU 现场信息}：寄存器值（通用寄存器、PC、PSW、栈指针）、指向页表的指针。
\end{enumerate}

\Q{Q3-1.4 创建进程主要完成哪些工作？}
\bk 课上强调的三件必做之事：\textbf{分配 PID、分配 PCB、分配地址空间}。完整步骤：
\begin{enumerate}
    \item 给新进程分配一个唯一的\textbf{进程标识符 PID}，并分配一个 \textbf{PCB}（若 PCB 表已满则创建失败）；
    \item 为进程分配\textbf{地址空间}；
    \item 初始化 PCB：设置各字段的默认值（如状态为 New，优先级默认值）；
    \item 设置相应的队列指针，例如把新进程\textbf{加入就绪队列}；
    \item 创建或扩充其他数据结构，如记账信息。
\end{enumerate}
UNIX 用 \cd{fork} + \cd{exec} 两步完成；Windows 用 \cd{CreateProcess} 一步完成（需要 10 个参数，既创建进程又装入程序）。

\Q{Q3-1.5 进程的生命周期内都会经历哪些变化？怎样表示这些变化？}
变化过程：被创建 $\to$ 在就绪、运行、等待之间反复转换（可能被挂起、激活）$\to$ 终止 $\to$ 资源被回收。表示方法有三种角度：(1) PCB 中的\textbf{状态字段}（如 xv6 的 \cd{p->state}、Linux 的 \cd{task_struct->state}）；(2) PCB 在\textbf{不同队列之间移动}（就绪队列、各种等待队列）；(3) 画成\textbf{状态转换图}。

\Q{Q3-1.6 进程有哪些状态？进程状态之间的转换（条件？操作？）}
三种基本状态：
\begin{itemize}
    \item \textbf{运行态}（Running）：占有 CPU，正在 CPU 上运行。
    \item \textbf{就绪态}（Ready）：已经具备运行条件，只是暂时没有空闲的 CPU 而不能运行（“万事俱备，只欠 CPU”）。
    \item \textbf{等待态}（Waiting/Blocked，又称阻塞态、封锁态、睡眠态）：因等待某一事件（如 I/O 完成）而暂时不能运行，这时即使把 CPU 给它也没有用。
\end{itemize}
进程在消亡之前处于且仅处于其中一种状态。

\begin{btab}{P{0.152}P{0.424}P{0.424}}
\thead{转换 & 条件（触发事件） & 操作（由谁完成）}
\trow{就绪$\to$运行 & CPU 空闲，调度程序选中该进程 & 调度程序完成选择和上下文切换}
\trow{运行$\to$就绪 & \textbf{时间片用完}（时钟中断）；有更高优先级的进程就绪，当前进程被\textbf{抢占} & 中断处理后由调度程序切换}
\trow{运行$\to$等待 & 请求了暂时无法满足的事件：发出慢速 I/O 系统调用（如读盘）、\cd{wait} 等待子进程、申请不到锁、缺页需要读盘 & 进程\textbf{自己}执行阻塞原语（主动）}
\trow{等待$\to$就绪 & 所等待的事件发生：\textbf{I/O 完成中断}、子进程结束、锁被释放 & \textbf{其他实体}（中断处理程序或其他进程）执行唤醒原语（被动）}
\end{btab}
三状态之间理论上有 6 种转换，只保留 4 种：
\begin{itemize}
    \item \imp{没有“等待 $\to$ 运行”}：事件发生只表示进程具备了运行条件，它必须先进入就绪队列排队，由调度程序统一挑选。如果允许唤醒后直接运行，就绕过了调度程序，调度就失去了确定性和公平性。
    \item \imp{没有“就绪 $\to$ 等待”}：就绪进程并没有在执行，不可能发出“我要等待某个事件”的请求；只有正在运行的进程才能执行阻塞原语。
\end{itemize}
例：\cd{cat file} 的状态变化——运行（调用 \cd{read}）$\to$ 等待（等读盘）$\to$ 磁盘中断 $\to$ 就绪 $\to$ 被调度 $\to$ 运行 $\to$ 时间片用完 $\to$ 就绪 $\to$ ……

\Q{Q3-1.7 进程状态转换的发生，是否一定导致另一个转换发生？}
\textbf{不一定，但常常会}（以单 CPU 为例逐条分析）：
\begin{itemize}
    \item 运行 $\to$ 等待 / 终止：\textbf{会}。CPU 空出来了，如果就绪队列非空，必然有一个进程从就绪变为运行；终止还可能唤醒正在 \cd{wait} 的父进程（等待 $\to$ 就绪）。
    \item 运行 $\to$ 就绪：\textbf{会}。另一个就绪进程变为运行（如果就绪队列中只有它自己，就再次选中它自己）。
    \item 等待 $\to$ 就绪：\textbf{不一定}。在非抢占式调度下不会引起其他转换；在抢占式调度下，如果被唤醒的进程优先级比当前运行进程高，就会导致当前进程“运行 $\to$ 就绪”。
    \item 就绪 $\to$ 运行：一般是其他转换的\textbf{结果}，而不是原因。
\end{itemize}
就绪队列为空时，CPU 运行 idle 进程（xv6 中调度器执行 \cd{wfi} 等待中断）。

\Q{Q3-2.1 OS 给进程提供的内存空间，其地址是虚拟地址还是物理地址？为什么？}
\imp{是虚拟地址}（老师强调“很多学生容易混淆，需明确记住”）。

\textbf{证据：课件 myval 实验}。程序打印一个全局变量 \cd{myval} 的地址和值。关闭 ASLR 后同时运行 \cd{./myval 5} 和 \cd{./myval 6}，两个进程打印的地址都是 \cd{0x601044}，但值分别是 5 和 6。同一个物理地址不可能同时存放 5 和 6，所以这个地址只能是\textbf{两个不同地址空间中的同一个虚拟地址}，它们被映射到了不同的物理页。

\textbf{为什么要用虚拟地址}：(1) \textbf{隔离与保护}：每个进程只能访问自己页表中映射的内存，访问不到别的进程和内核；(2) \textbf{简化链接和加载}：所有程序都可以链接到相同的地址开始运行，不必关心实际被装到物理内存的哪里；(3) 提供\textbf{巨大、连续的地址空间}的幻觉：物理内存可以零散分配，甚至部分内容在磁盘上；(4) \textbf{方便共享}：共享库、共享内存、COW 都靠把不同进程的虚拟页映射到同一物理页实现。

\textbf{课件思考：为什么在现代系统上两次运行打印的地址不一样？} 因为 \textbf{ASLR}（地址空间布局随机化）：每次运行时栈、堆、共享库、以及 PIE 程序的代码段基址都随机变化，目的是防止缓冲区溢出攻击者预知地址（联系 CSAPP 第 3 章 attack lab）。关闭方法：\cd{sysctl -w kernel.randomize_va_space=0}，或用 \cd{setarch -R ./myval} 只对单个程序关闭；此外 gcc 默认生成 PIE 程序，加 \cd{-no-pie} 编译才会得到 \cd{0x60xxxx} 这样的固定地址。

\Q{Q3-2.2 OS 如何描述进程的地址空间？}
分两层描述：
\begin{itemize}
    \item \textbf{区域层（逻辑层）}：记录地址空间中\textbf{“应该有什么”}——每一段连续的区域的起止地址、访问权限、对应哪个文件的哪个偏移（还是匿名内存）。Linux 中，\cd{task_struct->mm} 指向 \cd{mm_struct}，其中维护一组 \textbf{VMA}（\cd{vm_area_struct}），每个 VMA 描述一段区域。VMA 组织成链表和树（AVL 树、红黑树，新版本为 maple tree，不同版本不同），树结构便于在 \cd{mmap} 时快速检查新区域是否与已有区域重叠，以及缺页时快速找到地址所属的区域。
    \item \textbf{页表层（硬件层）}：记录\textbf{“现在映射到了哪个物理页”}。\cd{mm_struct} 中的 \cd{pgd} 指向顶级页表，进程运行时它被装入 CR3（x86）或 \cd{satp}（RISC-V），供 MMU 使用。
\end{itemize}
两层的配合：发生缺页时，内核查 VMA 判断出错地址是否合法。如果落在某个 VMA 内，说明是合法访问，只是还没有调入，于是按需分配页、读入内容、填写页表；如果不在任何 VMA 内，就是非法访问，向进程发送 SIGSEGV（段错误）。xv6 的设计更简单：没有 VMA，只用 \cd{p->sz} 表示用户空间是从 0 到 \cd{sz} 的一段连续区域，再加上 \cd{p->pagetable}。

在 Linux 上可以用 \cd{cat /proc/PID/maps} 查看一个进程的所有 VMA（见 3.4）。

\Q{Q3-2.3 为什么有了进程后又引入线程？}
\begin{enumerate}
    \item \textbf{应用的需要}：一个应用中往往同时有多个活动，它们需要\textbf{并发}执行，又要\textbf{共享同一份数据}。例：字处理软件中，一个线程与用户交互、一个线程在后台重新排版、一个线程定时自动保存，三者操作的是同一份文档；Web 服务器同时处理多个请求，这些请求共享网页缓存。如果用多个进程来做，进程之间共享数据很麻烦（要用共享内存、管道等进程间通信机制）。
    \item \textbf{开销的考虑}：创建、撤销、切换进程都要建立或切换地址空间、页表，切换后 TLB 和 cache 失效；进程间通信要经过内核。线程共享所在进程的地址空间，创建只需分配栈和 TCB，同进程线程之间切换不用换页表，通信直接读写共享变量，代价小得多。进程开销大\textbf{限制了系统的并发度}。
    \item \textbf{性能的考虑}：在多核处理器上，一个进程的多个线程可以\textbf{真正并行}地运行在不同核上。
\end{enumerate}
一句话：线程把进程的两个属性——\textbf{资源的拥有者}和\textbf{CPU 调度的单位}——拆开了。

\Q{Q3-2.4 怎样实现线程机制？为什么有各种支持线程的方式？}
三种实现方式：\textbf{用户级线程}（N:1，在用户空间由线程库实现，内核不知道线程的存在）、\textbf{核心级线程}（1:1，由内核实现和调度）、\textbf{混合方式}（M:N，用户级线程多路复用到若干核心级线程上）。详见 3.6 节的对比表。

之所以有多种方式，是因为它们之间存在\textbf{不可兼得的权衡}：切换开销（用户级最小）与阻塞隔离、多核并行（核心级才能做到）不可兼得；调度策略能否由应用定制；对 OS 有没有要求（用户级可以在任何 OS 上实现）。另外还有历史原因：早期内核不支持线程，只能在用户态实现；后来内核支持线程了，又嫌内核级线程切换太慢；混合方式想兼得两者的优点，但实现非常复杂。今天 Linux（NPTL）和 Windows 都采用 1:1；而 Go 的 goroutine、Java 的虚拟线程又是 M:N——这是概念的回归。

\Q{Q3-2.5 线程包 Pthreads 中相关函数的功能？}
\begin{itemize}
    \item \cd{pthread_create}：创建一个新线程，指定它要执行的函数和传给该函数的参数，返回新线程的标识符。
    \item \cd{pthread_exit}：结束调用它的线程，并给出返回值。
    \item \cd{pthread_join}：等待某个特定线程退出，并回收它的资源，类似进程的 \cd{wait}。注意：如果 \cd{main} 不等待其他线程就返回，整个进程（包括所有线程）都会退出。
    \item \cd{pthread_yield}（POSIX 标准名为 \cd{sched_yield}）：调用线程\textbf{自愿}释放 CPU，让其他线程运行。
    \item \cd{pthread_attr_init} / \cd{pthread_attr_destroy}：创建并初始化 / 删除一个线程属性结构（栈大小、调度策略、优先级等）。
    \item 同步相关：\cd{pthread_mutex_lock/unlock}（互斥锁），\cd{pthread_cond_wait/signal}（条件变量）。
\end{itemize}
\textbf{课件强调“线程是自愿让出 CPU 的”}，这是针对\textbf{用户级线程}说的：内核不知道线程的存在，时钟中断只能在进程之间切换，线程库无法打断一个正在运行的线程。所以线程必须主动调用 \cd{pthread_yield}，或者调用线程库提供的会阻塞的函数（如等锁），线程库才有机会切换到别的线程。如果一个线程陷入死循环，同一进程的其他线程就永远得不到运行。

\Q{Q3-3.1 中断/异常机制与进程线程模型的关联？}
\imp{进程状态的每一次转换，背后几乎都是一次中断、异常或系统调用的处理}——因为只有进入内核，OS 才有机会修改进程状态、切换进程。
\begin{itemize}
    \item 运行 $\to$ 就绪：\textbf{时钟中断}发现时间片用完；或 I/O 中断唤醒了更高优先级的进程，当前进程被抢占。
    \item 运行 $\to$ 等待：\textbf{系统调用}（\cd{read}、\cd{wait}、\cd{sleep}）或\textbf{缺页异常}（需要读盘）。
    \item 等待 $\to$ 就绪：\textbf{I/O 中断}（磁盘读完了）；或另一个进程的系统调用（如 \cd{exit} 唤醒父进程、释放锁）。
    \item 运行 $\to$ 终止：\cd{exit} \textbf{系统调用}；或发生异常（非法访问、除零）被 OS 杀死。
    \item 就绪 $\to$ 运行：发生在上述某次中断/异常/系统调用处理完毕、\textbf{返回用户态之前}的调度中。
\end{itemize}
另外，中断机制中的\textbf{保存现场}正是进程“被打断后还能恢复执行”的基础。

xv6 中一次时钟中断引起进程切换的完整链条：时钟中断 $\to$ \cd{uservec} 把用户寄存器保存到 trapframe $\to$ \cd{usertrap()} $\to$ \cd{devintr()} 返回 2 $\to$ \cd{yield()} 把状态改为 RUNNABLE $\to$ \cd{sched()} $\to$ \cd{swtch()} 切换到调度器线程 $\to$ \cd{scheduler()} 选出下一个进程 $\to$ \cd{swtch()} 切换过去 $\to$ 该进程从它自己上一次的 \cd{yield()} 中返回 $\to$ \cd{usertrapret()} $\to$ \cd{sret} 回到它的用户态。中断是进程模型运转的“引擎”。

\Q{Q3-3.2 机制与策略分离的原则在进程线程模型中的体现？}
\begin{itemize}
    \item 上下文切换、PCB、各种队列是\textbf{机制}；调度算法、时间片长度是\textbf{策略}。
    \item \cd{fork} 和 \cd{exec} 是机制；在 fork 和 exec 之间如何设置子进程（重定向、管道、改工作目录）是策略，由 shell 决定。\textbf{fork + exec 分离是最典型的例子}。
    \item 阻塞/唤醒原语是机制；唤醒的顺序（唤醒一个还是全部、按什么顺序）是策略。
    \item 内核提供的核心级线程、调度器激活（scheduler activation）是机制；用户级线程库采用什么调度算法是策略。
    \item 优先级字段和 \cd{nice} 系统调用是机制；谁更重要是策略，由用户决定。
    \item Linux 的 \cd{sched_class} 让调度算法以“调度类”的形式可插拔。
\end{itemize}

\Q{Q3-3.3 协程是什么？为什么引入协程？协程怎么用？}
\textbf{是什么}：协程（coroutine）是\textbf{用户态的、协作式（非抢占）}的执行单元。多个协程运行在\textbf{同一个线程}中，只在程序显式指定的让出点（\cd{yield}、\cd{await}）才切换，切换不经过内核，只需保存/恢复少量寄存器，比线程更轻量。它\textbf{与 OS 无关}，由编程语言或库实现，内核看到的只是普通线程。

\textbf{为什么引入（C10K 问题）}：服务器要同时处理上万个连接，而这些连接大部分时间在等网络 I/O。每个连接一个进程，开销太大；每个连接一个线程，每个线程要占一个栈（默认 8MB 虚拟空间）还要内核调度，上万个线程也扛不住，于是有了线程池；用 \cd{epoll} + 非阻塞 I/O + 回调函数（事件驱动、有限状态机）效率很高，但代码被拆散到各个回调中，难写难读（“回调地狱”）。协程的底层仍然是 \cd{epoll}，但让程序员能用\textbf{同步的写法获得异步的性能}：遇到 I/O 时 \cd{await}，协程挂起，事件循环去运行其他协程，I/O 就绪后再恢复它。

\textbf{怎么用}：Python 用 \cd{async def} 定义协程，在协程中用 \cd{await asyncio.sleep(1)} 让出，用 \cd{asyncio.gather()} 并发运行多个协程——3 个分别等待 1s、2s、3s 的任务总耗时约 3s 而不是 6s。Go 用 \cd{go f()} 启动 goroutine；Rust、C++20、JavaScript 也都有 async/await。

\textbf{缺点}与用户级线程相同：如果某个协程调用了阻塞的系统调用，或者进行长时间计算而不 \cd{await}，同一线程上的所有协程都会被卡住。\textbf{分类}：有栈协程（如 goroutine，每个协程有自己的栈，可以在任意函数深处挂起）与无栈协程（如 async/await，编译器把函数改写成状态机，只能在 \cd{await} 处挂起）。Windows 中类似的概念叫纤程（Fiber）。

\subsection{基本概念}

\textbf{顺序环境}：系统中只有一个程序在运行，它独占所有资源，执行不受外界影响。特征：(1) 程序执行的\textbf{顺序性}；(2) \textbf{封闭性}：独占资源，执行过程不受外界影响；(3) \textbf{结果可再现}：运行结果与执行速度无关。

\textbf{多道程序设计}：允许多个程序\textbf{同时}进入内存并运行，目的是提高系统效率。

\textbf{并发环境}：在一段时间间隔内，单处理器上有两个或两个以上的程序同时处于“已开始运行但尚未结束”的状态，并且它们执行的先后次序不是事先确定的。在并发环境中执行的程序称为并发程序。并发环境的特征：
\begin{itemize}
    \item \textbf{间断性}：程序“执行—停—执行”，走走停停；
    \item \textbf{资源共享}：多个程序共享系统资源；
    \item \textbf{独立性和制约性}：各程序独立推进，又因共享资源、相互协作而相互制约；
    \item \textbf{程序和计算不再一一对应}：同一段程序可以同时被多次执行，形成多个计算；
    \item \textbf{结果不可再现}：结果与各程序推进的相对速度有关。
\end{itemize}
在这样的环境中，静态的“程序”已不足以描述正在发生的事情，于是引入了动态的“进程”。

\textbf{进程的定义}（课件标色部分）：进程是具有独立功能的程序\textbf{关于某个数据集合上的一次运行活动}，是系统进行\textbf{资源分配和调度的独立单位}。进程又称任务（task）或作业（job）。理解要点：“一次运行活动”说明同一个程序运行两次就是两个进程；“数据集合”说明同一程序处理不同数据也是不同的进程。系统资源以进程为单位分配（内存、文件……），每个进程有独立的地址空间，OS 把 CPU 调度给进程。

\ez 引入线程以后：\imp{进程是资源分配的单位，线程是 CPU 调度的单位}。这种分工只在进程和线程同时存在时才成立；在只有线程、没有进程概念的系统中，线程就等同于传统的进程。

\textbf{进程与程序的区别}：(1) 进程能更准确地刻画并发，程序不能；(2) 程序是静态的，进程是动态的；(3) 进程有生命周期，有诞生有消亡，是短暂的，程序是相对长久的；(4) 一个程序可以对应多个进程；(5) 进程能创建其他进程，程序不能。比喻：程序是菜谱，进程是按菜谱做菜的过程——同一份菜谱可以同时有好几个厨师在做。

\textbf{进程的分类}：
\begin{itemize}
    \item \textbf{系统进程与用户进程}：系统进程完成 OS 的部分功能，如\textbf{守护进程}（daemon），平时睡眠，有任务时才被唤醒（打印守护进程、\cd{sshd}）。
    \item \textbf{前台进程与后台进程}：前台进程与用户交互；后台进程不占用终端。
    \item \textbf{CPU 密集型与 I/O 密集型}：CPU 密集型进程大部分时间在计算（如科学计算、视频编码）；I/O 密集型进程大部分时间在等 I/O，每次只用一小段 CPU（如编辑器、Web 服务器）。这个区分对调度很重要（第 4 章）。
\end{itemize}
\textbf{进程之间的关系}：UNIX 中有\textbf{进程家族树}，以 init（PID 1）为根，父进程创建子进程；父进程先结束时，子进程成为孤儿，由 init 收养。Windows 中进程地位平等，没有层次关系，创建者得到子进程的句柄，并且句柄可以转交给别的进程。

\textbf{如何查看系统中有多少进程}（课件问题）：Linux 下用 \cd{ps aux}（每行一个进程）、\cd{ls /proc} 中的数字目录（每个数字是一个 PID）、\cd{pstree} 查看家族树、\cd{top} 动态查看；xv6 中按 Ctrl-P 打印进程表。

\subsection{进程状态模型}

\begin{center}
\begin{tikzpicture}[>=Stealth, every node/.style={font=\scriptsize},
  st/.style={draw, rounded corners=2pt, inner sep=2pt, font=\footnotesize}]
\node[st] (new) at (0,0) {新建};
\node[st] (ready) at (2,0) {就绪};
\node[st] (run) at (4.6,0) {运行};
\node[st] (exit) at (6.6,0) {终止};
\node[st] (blk) at (3.3,-1.3) {阻塞};
\draw[->] (new) -- node[above]{接纳} (ready);
\draw[->] (ready) to[bend left=25] node[above]{调度选中} (run);
\draw[->] (run) to[bend left=25] node[below]{时间片到/抢占} (ready);
\draw[->] (run) -- node[above]{释放} (exit);
\draw[->] (run) -- node[right]{等待事件} (blk);
\draw[->] (blk) -- node[left]{事件发生} (ready);
\end{tikzpicture}
\end{center}

\textbf{进程的其他状态}：
\begin{itemize}
    \item \textbf{创建态}（New）：已经完成创建进程所必需的工作（分配了 PID、PCB），但\textbf{尚未同意执行}该进程，原因是资源有限。它的意义是在系统负载过重时起缓冲作用——老师的比喻是故宫门票预约：先登记，限流放行。\textbf{并非所有系统都有创建态}。
    \item \textbf{终止态}（Terminated）：进程终止执行后进入该状态，此时还可以完成一些数据统计工作（父进程读取退出码、记账），然后再回收资源。\textbf{主流系统都有终止态}（如 UNIX 的僵尸态）。
    \item \textbf{挂起态}（Suspend）：把一个进程从内存转到磁盘上，进程不再占用内存空间，其映像被交换到磁盘。用于调节系统负载，或者用户想暂停某个进程。
\end{itemize}

\textbf{三种模型的异同}（课件思考题）：
\begin{itemize}
    \item \textbf{三状态模型}只描述 CPU 的分配：运行、就绪、等待。
    \item \textbf{五状态模型}加上了创建态和终止态，描述了进程的出生和死亡，是\textbf{最接近实际系统}的模型。
    \item \textbf{七状态模型}再加上\textbf{就绪挂起}和\textbf{阻塞挂起}，引入了“是否在内存中”这一维度，它与“是否就绪”是正交的。新增的转换有：阻塞 $\to$ 阻塞挂起（内存紧张时把阻塞进程换出）、阻塞挂起 $\to$ 就绪挂起（等待的事件在换出期间发生了）、就绪挂起 $\to$ 就绪（激活，换入内存）、就绪 $\to$ 就绪挂起等。七状态模型在现代系统中很少使用，因为虚拟内存以页为单位换出，不需要把整个进程换出去。
\end{itemize}

\textbf{七状态进程模型}（Stallings）：\par
\begin{center}
\begin{tikzpicture}[>=Stealth, every node/.style={font=\scriptsize},
  st/.style={draw, rounded corners=2pt, inner sep=2pt, minimum width=1.3cm, minimum height=0.55cm, align=center, font=\footnotesize}]
\node[st] (new)  at (3.6,0)    {创建};
\node[st] (exit) at (7.4,0)    {退出};
\node[st] (rs)   at (0,-1.5)   {就绪挂起};
\node[st] (r)    at (3.6,-1.5) {就绪};
\node[st] (run)  at (7.4,-1.5) {运行};
\node[st] (bs)   at (0,-3)     {阻塞挂起};
\node[st] (b)    at (3.6,-3)   {阻塞};
\draw[->] (new) -- node[right]{提交} (r);
\draw[->] (new.west) to[bend right=15] node[above left=-1pt]{提交} (rs.north);
\draw[->] (rs.10) -- node[above]{激活} (r.170);
\draw[->] (r.190) -- node[below]{挂起} (rs.350);
\draw[->] (r.10) -- node[above]{调度} (run.170);
\draw[->] (run.190) -- node[below]{超时} (r.350);
\draw[->] (run.north west) .. controls (5.6,-0.55) and (1.2,-0.55) .. node[pos=0.25, above]{挂起} (rs.north east);
\draw[->] (run) -- node[right]{释放} (exit);
\draw[->] (run.south) |- node[pos=0.3, right]{等待事件} (b.east);
\draw[->] (b) -- node[right]{事件出现} (r);
\draw[->] (bs) -- node[right]{事件出现} (rs);
\draw[->] (bs.10) -- node[above]{激活} (b.170);
\draw[->] (b.190) -- node[below]{挂起} (bs.350);
\end{tikzpicture}
\end{center}
读图要点：横向看是“是否在内存中”（左列挂起在磁盘，中列在内存），纵向看是“是否在等事件”（上行就绪，下行阻塞）。挂起只改变前一个维度：\textbf{阻塞挂起的进程等到事件后变为就绪挂起}，仍在磁盘上，要\textbf{激活}换入内存后才能被调度。

\textbf{Linux 进程状态}（课件第 16 页，共 6 个方框）：\par
\begin{center}
\begin{tikzpicture}[>=Stealth, every node/.style={font=\scriptsize},
  st/.style={draw, rounded corners=2pt, inner sep=2pt, align=center, font=\scriptsize}]
\node[st] (rdy) at (4,0)     {\texttt{TASK\_RUNNING}\\就绪};
\node[st] (cpu) at (4,-2.2)  {占有 CPU 执行\\（也是 \texttt{TASK\_RUNNING}）};
\node[st] (un)  at (0.6,-1.1) {\texttt{TASK\_UNINTERRUPTIBLE}\\深度睡眠（D）};
\node[st] (in)  at (7.6,-1.1) {\texttt{TASK\_INTERRUPTIBLE}\\浅度睡眠（S）};
\node[st] (stp) at (0.6,-3.5) {\texttt{TASK\_STOPPED}\\暂停（T）};
\node[st] (zom) at (7.6,-3.5) {\texttt{TASK\_ZOMBIE}\\死亡但 PCB 未注销（Z）};
\draw[<-] (rdy.north) -- ++(0,0.45) node[above]{\texttt{fork()}};
\draw[->] (rdy.250) -- node[left]{\texttt{schedule()}} (cpu.110);
\draw[->] (cpu.70) -- node[right]{时间片到} (rdy.290);
\draw[->] (cpu.west) -| node[pos=0.25, below, align=center]{等待资源\\\texttt{sleep\_on()}} (un.south);
\draw[->] (un.north) |- node[pos=0.6, above]{资源到位 \texttt{wake\_up()}} (rdy.180);
\draw[->] (cpu.east) -| node[pos=0.25, below, align=center]{等待资源\\\texttt{interruptible\_sleep\_on()}} (in.south);
\draw[->] (in.north) |- node[pos=0.6, above, align=center]{资源到位或收到信号\\\texttt{wake\_up\_interruptible()}} (rdy.0);
\draw[->] (cpu.230) |- node[pos=0.7, above]{\texttt{ptrace()}、SIGSTOP} (stp.east);
\draw[->] (stp.west) -- ++(-0.9,0) |- node[pos=0.1, right, align=left]{SIGCONT\\\texttt{wake\_up()}} ($(rdy.north)+(-0.6,0.8)$) -- ([xshift=-0.6cm]rdy.north);
\draw[->] (cpu.310) |- node[pos=0.7, above]{\texttt{do\_exit()}} (zom.west);
\end{tikzpicture}
\end{center}
与五状态模型对比：(1) Linux 用\textbf{同一个} \cd{TASK_RUNNING} 表示就绪和运行，两者只靠“是否正在某个 CPU 上”区分；(2) 阻塞态细分为\textbf{可被信号唤醒}的浅睡与\textbf{不响应信号}的深睡；(3) 多了供调试和作业控制用的暂停态；(4) 终止态对应 ZOMBIE；(5) 没有显式的创建态和挂起态（换出以页为单位，不挂起整个进程）。

\textbf{实际系统的进程状态}：\par
\begin{btab}{P{0.130}P{0.870}}
\trow{Linux & \cd{TASK_RUNNING}：\imp{同时包含就绪和运行}两种状态（\cd{ps} 中显示为 R）。\cd{TASK_INTERRUPTIBLE}：浅度睡眠，可以被信号唤醒（S）。\cd{TASK_UNINTERRUPTIBLE}：深度睡眠，不响应信号（D），即使 \cd{kill -9} 也杀不掉，典型场景是等待磁盘 I/O——如果中途被信号打断，内核数据结构很难恢复到一致状态。\cd{TASK_STOPPED}：暂停（T），收到 SIGSTOP、按 Ctrl-Z 或被 ptrace 调试时进入。\cd{EXIT_ZOMBIE}：僵尸（Z），进程已经结束，但父进程还没有调用 \cd{wait}，PCB 保留着退出码。转换：\cd{fork} 创建后为 RUNNING；\cd{wake_up} 把睡眠进程改为 RUNNING；\cd{do_exit} 进入 ZOMBIE}
\trow{xv6 & \cd{enum procstate {UNUSED, USED, SLEEPING, RUNNABLE, RUNNING, ZOMBIE}}。UNUSED：空闲的 PCB 槽位；USED：刚被 \cd{allocproc} 分配，正在初始化（相当于创建态）；RUNNABLE：就绪；RUNNING：运行；SLEEPING：阻塞，\cd{p->chan} 记录在等什么；ZOMBIE：已退出，等待父进程 \cd{wait}，被 \cd{freeproc} 回收后回到 UNUSED。UNUSED/USED 表示的是 PCB 槽位的使用情况，不算常规的进程状态。xv6 没有挂起状态}
\end{btab}

\ez \textbf{用户级线程的状态与进程状态是两套}：线程执行阻塞系统调用时，内核阻塞的是\textbf{整个进程}，但线程库中这个线程的状态\textbf{仍然显示为 running}（实际上并没有在运行）。Stallings 的三个例子（进程 B 有线程 1、2，线程 2 正在运行）：
\begin{enumerate}
    \item 线程 2 执行阻塞的 I/O 系统调用：内核把进程 B 置为阻塞；线程库中线程 2 仍是“运行”。
    \item 时钟中断，进程 B 时间片用完：进程 B 变为就绪；线程 2 仍是“运行”。
    \item 线程 2 调用线程库的同步函数等待线程 1：线程库把线程 2 置为阻塞、切换到线程 1 运行；在内核看来进程 B 一直在运行。
\end{enumerate}

\subsection{PCB、地址空间与进程队列}

\textbf{PCB}：PCB 是进程存在的\textbf{唯一标志}，与进程一一对应。所有进程的 PCB 集合称为\textbf{进程表}，它在内存中\textbf{预先分配}好（大小决定了系统最多能有多少进程）。实例：Linux 的 \cd{task_struct}、xv6 的 \cd{struct proc}、Windows 的 EPROCESS/KPROCESS/PEB（分别位于执行体、内核和用户空间）、Solaris 的 \cd{proc_t}。

\textbf{课件强调的 PCB 核心内容}：(1) \textbf{可执行文件的位置}——能通过它（inode）找到磁盘上的程序文件，缺页时才能从文件中读入代码；(2) \textbf{地址空间}——\cd{mm} 指针；(3) \textbf{打开文件表}——子进程继承父进程的打开文件，默认打开 0（标准输入）、1（标准输出）、2（标准错误）；以及 (4) \textbf{上下文}，其中 PC、PSW、栈指针、\textbf{页表指针}最重要。为什么要在 PCB 中保存现场？因为进程被暂停后还要从断点继续执行，必须能恢复它被打断时的 CPU 状态。

\textbf{Linux \cd{task_struct} 的主要字段}：
\begin{itemize}
    \item 标识：\cd{pid}（实际上是线程 ID）和 \cd{tgid}（线程组 ID，即用户看到的“进程 ID”，\cd{getpid()} 返回的是它）。
    \item 状态 \cd{state}；亲缘关系 \cd{real_parent}、\cd{parent}、\cd{children}、\cd{sibling}（构成家族树）。
    \item 调度相关：\cd{prio}、\cd{static_prio}、\cd{normal_prio}、\cd{rt_priority}、\cd{sched_class}、\cd{policy}。
    \item 统计：\cd{utime}/\cd{stime}（用户态/内核态运行时间）、\cd{nvcsw}/\cd{nivcsw}（自愿/非自愿上下文切换次数）。
    \item 权限 \cd{uid}、\cd{euid}、\cd{cred}；信号 \cd{signal}、\cd{sighand}、\cd{blocked}、\cd{pending}。
    \item 资源：\textbf{\cd{mm}}（地址空间）、\textbf{\cd{fs}}（根目录、当前目录）、\textbf{\cd{files}}（打开文件表）；\cd{thread}（保存的寄存器）、\cd{stack}（内核栈）。
\end{itemize}
资源都是通过\textbf{指针}引用的：\cd{fork} 时复制这些结构；创建线程时则让新任务直接共享这些指针（\cd{clone} 时带 \cd{CLONE_VM}、\cd{CLONE_FILES} 等标志）。所以 \textbf{Linux 中进程与线程的区别仅在于与创建者共享多少资源}。

\textbf{xv6 \cd{struct proc}}：需要持有 \cd{p->lock} 才能访问的字段有 \cd{state}、\cd{chan}（睡眠在哪个通道上）、\cd{killed}、\cd{xstate}（退出状态）、\cd{pid}；\cd{parent} 由全局的 \cd{wait_lock} 保护；进程私有的字段有 \cd{kstack}（内核栈地址）、\cd{sz}（用户内存大小）、\cd{pagetable}、\cd{trapframe}（用户态寄存器）、\cd{context}（内核态寄存器，供 \cd{swtch} 使用）、\cd{ofile[NOFILE]}（打开文件）、\cd{cwd}（当前目录）、\cd{name}。按四类划分：描述信息（\cd{pid}、\cd{name}、\cd{parent}）、控制信息（\cd{state}、\cd{chan}、\cd{killed}、\cd{xstate}）、资源（\cd{pagetable}、\cd{sz}、\cd{ofile}、\cd{cwd}）、现场（\cd{trapframe}、\cd{context}）。xv6 没有优先级、用户 ID、信号和线程。

\textbf{进程地址空间}：从高地址到低地址依次是：内核空间 | 用户栈（向下增长）| 共享库和 \cd{mmap} 区域 | 堆（向上增长）| 数据段（\cd{.data}、\cd{.bss}）| 代码段。虚拟地址空间有四个特点：\textbf{巨大、布局统一、安全、稀疏}。

谈地址空间“有多大”时，要\textbf{先说清场景}：x86-64 Linux 使用 4 级页表，虚拟地址有效位为 48 位（$4\times9+12$），用户空间 128TB；RISC-V Sv39 为 39 位，xv6 的 \cd{MAXVA} 实际只用到 $2^{38}$（避免符号扩展问题）。

“稀疏”意味着大部分地址并没有使用，所以不能用一个大数组描述整个地址空间，而是只为\textbf{有内容的区域}建立 VMA，记录起止地址、权限、标志和对应的文件 \cd{vm_file}。

\textbf{查看 \cd{/proc/PID/maps}}：每一行就是一个 VMA。权限为 \cd{r-xp} 的是代码段，\cd{r--p} 是只读数据，\cd{rw-p} 是可写数据；\cd{[heap]} 是堆（程序调用过 \cd{malloc} 才会出现）；libc 和动态链接器 ld 各有几段（\cd{---p} 是不可访问的保护空洞）；\cd{[stack]} 是主线程栈；\cd{[vdso]} 是内核映射给用户的代码。权限最后一位 \cd{p} 表示私有映射（写时复制），\cd{s} 表示共享映射；有文件路径的是文件映射，设备号和 inode 为 \cd{00:00 0} 的是匿名映射。编写观察用的测试程序时，要让它用到堆和只读常量，并且不能立刻退出（如在末尾调用 \cd{getchar()} 等待输入），否则来不及查看。

创建进程时\textbf{并不真正读入}代码和数据，只建立 VMA 映射；堆和栈的内容也在运行时动态产生（与导论中“第一条指令就缺页”相呼应）。

\textbf{进程队列}：OS 为每一类进程建立一个或多个队列，队列元素是 PCB；进程状态改变，就是把它的 PCB \textbf{从一个队列移到另一个队列}。等待不同事件的进程可以排在不同的等待队列中；就绪队列可以按优先级分成多个；单 CPU 时，“运行队列”中只有一个进程。xv6 没有显式的队列：调度器线性扫描 \cd{proc[]} 数组找 RUNNABLE 的进程；“等待队列”用 \cd{p->chan} 字段隐式表示，\cd{wakeup(chan)} 扫描整个进程表，唤醒所有睡在该通道上的进程（时间复杂度 O(n)）。

\subsection{进程控制}

\textbf{进程控制}的主要操作有：创建、撤销（终止）、阻塞、唤醒、挂起、激活、改变优先级等。它们都由\textbf{原语}完成。

\textbf{原语}（primitive）：完成某种特定功能的一段程序，具有\textbf{不可分割性}（不可中断性），即执行过程必须是连续的，不允许被打断。最基本的实现方法是在单 CPU 上\textbf{关中断/开中断}；多 CPU 上关中断不够，还要加锁（xv6 用 \cd{p->lock}）。\textbf{为什么必须原子执行}：以阻塞原语为例，它要“修改状态 + 把 PCB 移入等待队列 + 调用调度程序”。如果做到一半被打断——状态已经改为阻塞，但还没有挂入任何等待队列——这个进程就“丢失”了，以后再也没有人能唤醒它。

\textbf{进程的创建}：何时创建进程？系统初始化时（创建 init、各种守护进程）；提供 OS 服务时；用户登录时；正在运行的进程执行了创建进程的系统调用（\cd{fork}）；用户在命令行提交程序时。创建的步骤见 Q3-1.4。

\textbf{进程的终止}：
\begin{itemize}
    \item 自愿：正常退出、出错退出（如找不到文件后调用 \cd{exit(1)}）。
    \item 非自愿：严重错误（执行非法指令、段错误、除零、在用户态执行特权指令）；被其他进程杀死（\cd{kill}）或父进程终止；OS 干预（超出时限、内存不足、I/O 失败）。
\end{itemize}
\textbf{进程的撤销}：先结束它的子进程或线程，再回收资源（关闭打开的文件、断开网络连接、回收分配的内存），最后撤销 PCB。UNIX 中撤销分为\textbf{两个阶段}：进程调用 \cd{exit} 后，释放大部分资源，进入僵尸状态，PCB 仍保留着退出码；父进程调用 \cd{wait} 读走退出状态后，PCB 才被真正释放。如果父进程一直不 \cd{wait}，僵尸进程就会越积越多，耗尽 PCB。

\textbf{进程的阻塞与唤醒}：处于运行态的进程在执行过程中期待某一事件发生（如等待键盘输入、等待磁盘传输完成、等待其他进程发消息），当被等待的事件未发生时，\textbf{由进程自己执行阻塞原语}，使自己由运行态变为阻塞态——阻塞是\textbf{主动}的。唤醒则是\textbf{被动}的，由事件的发生者执行唤醒原语（如磁盘中断处理程序唤醒等待读盘的进程）。UNIX 中相关的调用有 \cd{wait}、\cd{sleep}、\cd{pause}、\cd{read}；Windows 中有 \cd{WaitForSingleObject}。

xv6 的实现：\cd{sleep(chan, lk)} 先获取 \cd{p->lock}，再释放调用者持有的锁 \cd{lk}，然后设置 \cd{p->chan = chan; p->state = SLEEPING;} 并调用 \cd{sched()} 切换走；被唤醒后重新获取 \cd{lk}。\cd{wakeup(chan)} 扫描整个进程表，把所有 SLEEPING 且 \cd{chan} 相同的进程改为 RUNNABLE。

\textbf{UNIX 的四个进程控制系统调用}：
\begin{itemize}
    \item \cd{fork}：通过\textbf{复制调用进程}来建立新进程，是最基本的进程创建方式。
    \item \cd{exec}：一族系统调用（\cd{execl}、\cd{execvp}……），用一段新的程序代码\textbf{覆盖}原来的地址空间，实现进程执行代码的转换。
    \item \cd{wait}：提供初级的进程同步，使一个进程等待另一个（子）进程结束。
    \item \cd{exit}：终止一个进程的运行。
\end{itemize}
OSTEP 第 5 章的例子 \cd{p4.c}：子进程先 \cd{close(1)} 关闭标准输出，再 \cd{open("p4.output", ...)}（新文件会得到最小的可用描述符，即 1），然后 \cd{execvp("wc", ...)}；父进程 \cd{wait} 等待子进程结束。效果等同于 shell 命令 \cd{wc p4.c > p4.output}。

\textbf{UNIX \cd{fork} 的实现步骤}（课件）：(1) 为子进程分配一个空闲的 proc 结构（PCB）；(2) 赋予子进程唯一的 PID；(3) \textbf{以一次一页的方式复制父进程的地址空间}；(4) 子进程继承父进程共享的资源，如打开的文件和当前工作目录；(5) 把子进程状态设为就绪，插入就绪队列；(6) 对子进程返回 \textbf{0}；(7) 对父进程返回\textbf{子进程的 PID}。

\textbf{xv6 \cd{fork()} 的实现}：调用 \cd{allocproc()}（找到一个 UNUSED 槽位、分配 PID、分配 trapframe 页、创建空的用户页表、把 \cd{context.ra} 设为 \cd{forkret}、\cd{context.sp} 设为内核栈顶）$\to$ \cd{uvmcopy()} 把父进程的用户内存\imp{逐页真正复制}一份（\textbf{上游 xv6 没有 COW}，COW 是实验课的一个 lab）$\to$ \cd{*np->trapframe = *p->trapframe} 复制用户寄存器，并令 \cd{np->trapframe->a0 = 0}（子进程的返回值）$\to$ 对每个打开文件调用 \cd{filedup} 增加引用计数，\cd{idup} 当前目录 $\to$ 设置 \cd{np->parent = p} $\to$ 状态设为 RUNNABLE $\to$ 父进程返回子进程 PID。

\textbf{“一次调用，两次返回”的真相}：子进程\textbf{从来没有调用过 fork}，它是被内核“造”出来的——它的 trapframe 是父进程的副本，只是 \cd{a0} 被改成了 0；它的内核上下文中 \cd{ra} 指向 \cd{forkret}。子进程第一次被调度时，\cd{swtch} 返回到 \cd{forkret}，再经 \cd{usertrapret} 和 \cd{sret} 回到用户态，从 trapframe 中恢复的 PC 正是父进程 \cd{ecall} 的下一条指令，\cd{a0} 为 0，于是看起来就像 fork “返回了 0”。

\textbf{写时复制（COW）} \bk：
\begin{enumerate}
    \item \cd{fork} 时只复制 \cd{task_struct}、\cd{mm_struct}、VMA 和\textbf{页表}，\textbf{不复制物理页}。父子进程的页表项指向相同的物理页，并且都被标记为\textbf{只读}，VMA 标记为私有（可写时复制）。
    \item 之后无论父进程还是子进程对某一页执行\textbf{写}操作，硬件都会因为写只读页而产生\textbf{页错误（保护性异常）}。
    \item 内核的缺页处理程序查 VMA，发现该区域本来是可写的私有映射，于是判定这是 COW：分配一个新物理页，复制原页内容，修改出错进程的页表项指向新页并设为可写（若原页只剩一个使用者，则直接改为可写，不必复制）。
    \item 返回用户态，\textbf{重新执行那条写指令}（这是故障类异常），这次写入成功。
\end{enumerate}
\textbf{COW 省下了什么}：\cd{fork} 之后通常紧跟着 \cd{exec}，而 \cd{exec} 会直接丢弃旧的地址空间、为新程序建立新映射。如果 \cd{fork} 时真的复制了所有页，这些复制就白做了；有了 COW，只有真正被写的页才会被复制。\ez \imp{COW 不是 Linux 发明的}，它是存储管理中早已有的技术，Linux 只是拿来优化 fork（往年 40\%–50\% 的考生答错）。其他优化：\cd{vfork}（子进程直接借用父进程的地址空间，父进程暂停直到子进程 exec 或 exit）、\cd{posix_spawn}（一步完成创建和装入）。

\subsection{线程模型}

\textbf{为什么引入线程}：应用需要、开销小、多核性能（详见 Q3-2.3）。\textbf{开销对比}：创建/撤销线程只需分配/回收栈和 TCB，而进程要建立/回收整个地址空间；同进程的线程切换不需要换页表，TLB 不会失效；线程间通信直接读写共享变量，不经过内核。

\textbf{例：Web 服务器的三种构造方法}（Tanenbaum）：
\begin{enumerate}
    \item \textbf{多线程}：一个\textbf{分派线程}（dispatcher）从网络接收请求，把它交给一个空闲的\textbf{工作线程}（worker）；工作线程先查共享的网页缓存，不命中就读磁盘——读盘时只阻塞这个工作线程，其他线程继续服务别的请求。特点：有并行性、使用阻塞的系统调用，编程简单。
    \item \textbf{单线程进程}：一次只处理一个请求，读盘时整个服务器空等。特点：无并行性、阻塞的系统调用，最简单但最慢。
    \item \textbf{有限状态机}：使用非阻塞的 I/O。发起读盘后不等待，而是把当前请求的状态记录在表中，转去处理下一个事件；I/O 完成时（通过中断、信号或 \cd{epoll} 通知）再根据记录的状态继续处理。特点：有并行性、非阻塞系统调用、依赖中断，性能高但编程困难。这就是 CSAPP 第 12 章的 I/O 多路复用，Nginx、Redis、Node.js 都采用这种方式。
\end{enumerate}
三者的本质区别在于\textbf{系统调用是否阻塞}：多线程用“阻塞调用 + 多个执行流”获得并行，状态机用“非阻塞调用 + 一个执行流”获得并行。课件思考题“没有线程的话速度会怎样”：单线程方案最慢，状态机方案快但难写。

\textbf{线程的概念}：进程有两个基本属性：\textbf{资源的拥有者}（虚拟地址空间、打开的文件、I/O 设备、保护）和 \textbf{CPU 调度的单位}（执行路径、状态、优先级）。线程把后者独立出来：线程是进程中的一个运行实体，是 \textbf{CPU 的调度单位}，有时也称为轻量级进程（LWP）。
\begin{itemize}
    \item \textbf{线程私有的}：程序计数器 PC、寄存器（上下文）、\textbf{栈和栈指针}、线程状态，以及线程局部存储（TLS）。
    \item \textbf{同一进程中线程共享的}：地址空间（代码、全局变量、堆）、打开的文件、子进程、定时器、信号及信号处理程序、账户信息。
\end{itemize}
栈必须私有，因为每个线程有自己的函数调用链，栈帧记录着各自的局部变量和返回地址。\ez \textbf{课件思考：线程 1 能否访问线程 2 的栈？} \imp{能}。线程栈之间\textbf{没有内存保护}，它们都在同一个地址空间里；如果线程 1 拿到了指向线程 2 栈上变量的指针，就可以读写它。危险在于：如果线程 2 的函数已经返回，这个指针就悬空了，读到的是被覆盖的数据。

\begin{btab}{P{0.160}P{0.280}P{0.280}P{0.280}}
\thead{& 用户级 ULT（N:1） & 核心级 KLT（1:1） & 混合（M:N）}
\trow{谁管理 & 用户空间的线程库 + 运行时系统，每个进程有自己的线程表 & 内核，内核中有进程表和线程表 & 在用户空间创建线程，由核心级线程（LWP）承载，内核调度核心级线程}
\trow{内核是否感知线程 & 否，内核只知道进程 & 是，以线程为单位调度 & 感知 LWP / 核心级线程}
\trow{一个线程阻塞 & \imp{整个进程阻塞} & 只阻塞该线程 & 只阻塞对应的核心级线程，运行时可把其他用户线程移到别的核心级线程上}
\trow{能否多核并行 & \imp{不能}（内核只给进程分配一个 CPU） & 能 & 能（最多同时用 N 个核）}
\trow{切换开销 & 极小（不进内核，只需几条保存/恢复寄存器的指令，类似 \cd{setjmp/longjmp}） & 大（要进内核；但同进程线程切换不换地址空间，仍比进程切换快） & 介于两者之间}
\trow{调度策略 & 可以由应用定制 & 由内核决定 & 两层调度配合（调度器激活）}
\trow{对 OS 的要求 & 可在任何 OS 上实现 & 需要 OS 支持 & 实现复杂}
\trow{典型 & 早期的 Pthreads 实现、GNU Pth & Windows、今天的 Linux NPTL & 早期 Solaris、Go GMP、Java 虚拟线程}
\end{btab}
\textbf{推导的出发点只有一个：内核是否感知线程}。内核不知道线程（用户级）$\Rightarrow$ 内核调度的单位是进程 $\Rightarrow$ 一个线程阻塞时内核只能阻塞整个进程；一个进程同一时刻只能在一个 CPU 上运行 $\Rightarrow$ 不能多核并行；线程切换不需要进内核 $\Rightarrow$ 开销小。用户级线程的其他缺点：线程必须\textbf{自愿}让出 CPU（时钟中断交给的是内核，线程库无从抢占）；发生缺页时同样阻塞整个进程。

\textbf{用户级线程阻塞问题的两种补救}：
\begin{enumerate}
    \item \textbf{改为非阻塞的系统调用}（称为 jacketing / wrapper）：在真正发出系统调用之前，先用 \cd{select}/\cd{poll} 或以 \cd{O_NONBLOCK} 方式试探这次调用会不会阻塞；如果会，就不发出调用，而是把该线程标记为阻塞，切换到别的线程，稍后再重试。这种方法不需要修改内核。
    \item \textbf{重写 I/O 库函数}：线程库提供一整套自己的 \cd{read}、\cd{write}、\cd{accept}（如 GNU Pth 的 \cd{pth_read}）来替换 libc 中的版本，对应用程序透明，其内部通常就用方法 1 实现。
\end{enumerate}
局限：对普通磁盘文件，\cd{select} 总是报告“可读”，\cd{O_NONBLOCK} 也不起作用；缺页更是无法包装。这些情况只能依靠异步 I/O（如 \cd{io_uring}）或核心级线程来解决。

\ez \textbf{Pthreads 只是一套 API 标准}（POSIX 线程接口），既可以用用户级线程实现，也可以用核心级线程实现。课件把它当作用户级线程的例子，是因为早期的实现多在用户态；\imp{今天 Linux 的 NPTL 是 1:1 的核心级线程}：\cd{pthread_create} 内部调用 \cd{clone} 系统调用，并带上 \cd{CLONE_VM}、\cd{CLONE_THREAD} 等共享标志，用 \cd{ps -eLf} 可以看到每个线程有自己的 TID。答题时可以写“Pthreads 是 POSIX 线程接口，早期多为用户级实现”，避开这个坑。

\textbf{线程是否越多越好？} 不是。每个线程都要占用一个栈（Linux 默认 8MB 虚拟空间）和内核数据结构；线程越多，切换开销和 cache/TLB 污染越严重，锁竞争也越激烈。经验：CPU 密集型任务的线程数约等于核数；I/O 密集型任务可以多于核数。所以服务器常用\textbf{线程池}（预先创建固定数量的线程，反复使用）；高并发 I/O 场景则改用协程或事件驱动。

\textbf{可再入（reentrant）程序}：能被多个进程（线程）同时调用而不出错的程序。条件：(1) 是\textbf{纯代码}，执行过程中\textbf{不修改自身}；(2) 不依赖全局/静态变量保存状态，所需的数据区由调用者提供（状态放在栈上或参数中）。例：多个进程共享同一份 libc 代码。反例：\cd{strtok}、\cd{rand()} 依赖静态变量保存上一次的状态，两个线程同时调用会互相干扰，所以不可再入（可再入版本是 \cd{strtok_r}、\cd{rand_r}）。

\textbf{C10K 相关关键词的演进}：每个连接一个进程（切换开销大）$\to$ 每个连接一个线程 $\to$ 线程池 $\to$ I/O 多路复用（\cd{select}/\cd{poll}/\cd{epoll}）+ 非阻塞/异步 I/O $\to$ 事件驱动（回调）$\to$ 协程（用同步的写法写异步代码）。这条路线适用于 I/O 密集型的场景。

\subsection{xv6 源码阅读：shell、fork、exec、wait（api）}

\textbf{shell 的主循环}（\cd{user/sh.c} 的 \cd{main}）：
\begin{enumerate}
    \item \cd{getcmd()} 读入一行命令。
    \item 如果是 \cd{cd} 命令，由 shell \textbf{自己}调用 \cd{chdir}，不 fork。原因：如果在子进程中执行 \cd{chdir}，改变的只是子进程的当前目录，子进程退出后 shell 的目录不变——这也反过来证明了父子进程的资源是相互独立的。
    \item 其他命令：\cd{if(fork1() == 0) runcmd(parsecmd(buf));}，即子进程解析并执行命令。
    \item 父进程调用 \cd{wait(0)} 等待子进程结束，然后打印提示符，读下一条命令。
\end{enumerate}
\textbf{\cd{runcmd()} 按命令类型处理}：
\begin{itemize}
    \item EXEC：调用 \cd{exec(argv[0], argv)}，成功则不再返回（地址空间已被替换）。
    \item REDIR（重定向）：先 \cd{close(fd)} 关闭要重定向的描述符，再 \cd{open(file)} 打开文件（得到刚关闭的那个最小描述符），然后对子命令递归调用 \cd{runcmd}。
    \item PIPE（管道）：\cd{pipe(p)} 创建管道；fork 出左子进程，它 \cd{close(1); dup(p[1]);} 把标准输出接到管道写端，再执行左边的命令；fork 出右子进程，它 \cd{close(0); dup(p[0]);} 把标准输入接到管道读端，再执行右边的命令；父进程关闭管道两端，\cd{wait} 两次。
\end{itemize}

\textbf{以 \cd{echo hi > out.txt} 为例的完整过程}：sh fork 出子进程 $\to$ 子进程 \cd{close(1); open("out.txt")}，文件得到描述符 1 $\to$ \cd{exec("echo")}，子进程的内存被替换为 echo 的代码，但\imp{文件描述符表保持不变} $\to$ echo 调用 \texttt{write(1, "hi\textbackslash n", 3)}，写入的是文件 $\to$ \cd{exit(0)}，子进程变为 ZOMBIE，并唤醒父进程 $\to$ 父进程的 \cd{wait} 返回子进程 PID，子进程的 PCB 被回收（变为 UNUSED）$\to$ sh 打印提示符。

内核中对应的实现：\cd{exec.c} 读取 ELF 文件，建立新页表，把各个段装入内存，建立用户栈并压入参数，然后用新页表替换 \cd{p->pagetable}、释放旧的，设置 \cd{trapframe->epc = elf.entry}（入口地址）和 \cd{sp}——\textbf{完全不碰 \cd{p->ofile}}，所以打开的文件被保留。\cd{exit}：关闭所有打开的文件，把自己的子进程过继给 init，唤醒父进程，状态设为 ZOMBIE，调用 \cd{sched()} 切换走且永不返回。\cd{wait}：扫描进程表寻找自己的 ZOMBIE 子进程，找到则把退出状态 \cd{xstate} 拷贝给用户、调用 \cd{freeproc} 回收，返回其 PID；有子进程但都还没退出，则 \cd{sleep} 等待；没有子进程则返回 $-1$。

\textbf{为什么把创建进程分成 fork 和 exec 两步？好处是什么？}
\begin{enumerate}
    \item \textbf{在 fork 和 exec 之间可以修改子进程的环境}：I/O 重定向、建立管道、改变当前目录、设置环境变量、降低权限、设置信号处理方式等，而被执行的程序对此毫不知情。这使得程序和它的 I/O 目标解耦（echo 不需要知道自己的输出去了文件），是 UNIX“小工具 + 管道组合”哲学的基础。
    \item \textbf{接口简单、正交}：\cd{fork} 没有参数，\cd{exec} 只管装入程序；对比 Windows 的 \cd{CreateProcess} 需要 10 个参数来覆盖各种设置。
    \item \textbf{shell 自身得以保留}：如果 shell 直接 exec，它自己就被替换掉了。
    \item \textbf{fork 单独使用也有意义}：创建并发执行的副本，例如多进程服务器为每个连接 fork 一个子进程。
    \item 体现了\textbf{机制与策略分离}：内核提供 fork、exec 两个机制，如何组合（重定向、管道）由 shell 这个用户程序决定。
    \item 代价：fork 复制了整个地址空间，紧接着 exec 又把它丢掉，造成浪费——这正是 COW、\cd{vfork} 等优化的来由。
\end{enumerate}
补充：命令末尾加 \cd{&} 表示后台运行，shell 不调用 \cd{wait} 等待它结束。
